\chapter{Química computacional: Hartree--Fock e DFT}\label{ch:b3:computational-chemistry}

Acredita-se que o íon hidreto de hélio, \ce{HeH+}, tenha sido a primeira
molécula formada quando o Universo jovem esfriou: um átomo de hélio
segurando um próton. Ele foi produzido em um tubo de descarga de laboratório
em 1925, mas procurado no espaço sem sucesso durante décadas, até que sua
raia rotacional foi detectada em uma nebulosa planetária em 2019. Àquela
altura, seu comprimento de ligação, suas frequências vibracionais e os
comprimentos de onda de suas raias eram conhecidos havia muito tempo ---
por cálculo. Um programa de computador que resolve, de modo aproximado, a
\omterm{thm:b3:quantum-model-systems:schrodinger}{equação de Schrödinger} dos elétrons pode prever a forma, a energia e o
espectro de uma molécula que ninguém pôs em um frasco. Este capítulo explica
como funcionam esses programas, da separação entre núcleos e elétrons à
teoria de Hartree--Fock e à \omterm{def:b3:computational-chemistry:dft}{teoria do funcional da densidade}, e para que
seus números merecem ou não confiança.

\begin{recall}
O \cref{ch:b3:quantum-model-systems} deu o hamiltoniano, a \omterm{thm:b3:quantum-model-systems:schrodinger}{equação de
Schrödinger} e o átomo de hidrogênio; o \cref{ch:b3:many-electron-atoms}, o
\omterm{def:b3:many-electron-atoms:slater}{determinante de Slater} e a \omterm{def:b3:many-electron-atoms:exchange}{integral de troca}. O volume do 2.º ano construiu
orbitais moleculares como combinações lineares de orbitais atômicos (CLOA),
com integrais de sobreposição, de Coulomb e de ressonância e um determinante
secular, e resolveu o método de Hückel; usou também espectros de fotoelétrons
para ver as energias orbitais.
\end{recall}

\begin{omfigure}
\includegraphics[width=0.6\linewidth]{images/book4/ai/b3-03-cluster.jpg}

{\small Um cluster de computação: fileiras de servidores nos quais roda hoje
a maioria dos cálculos de química quântica, muitas vezes por horas ou dias por molécula.}
\end{omfigure}

\section{Born--Oppenheimer e a superfície de energia potencial}

O hamiltoniano de uma molécula contém as energias cinéticas de seus núcleos
e elétrons e todas as suas interações coulombianas. Os núcleos são pelo
menos 1836 vezes mais pesados que os elétrons: para as mesmas forças, eles
se movem muito mais devagar, e os elétrons se ajustam quase instantaneamente
a cada posição dos núcleos.

\begin{definition}[Aproximação de Born--Oppenheimer]\label{def:b3:computational-chemistry:born-oppenheimer}
Na \emph{aproximação de Born--Oppenheimer}\index{aproximacao de Born--Oppenheimer@aproximação de Born--Oppenheimer},
a \omterm{thm:b3:quantum-model-systems:schrodinger}{equação de Schrödinger} eletrônica é resolvida para núcleos
mantidos fixos nas posições $\mathbf R$,
\[
\hat H_{\mathrm{el}}\,\psi_{\mathrm{el}}(\mathbf r;\mathbf R) =
E_{\mathrm{el}}(\mathbf R)\,\psi_{\mathrm{el}}(\mathbf r;\mathbf R),
\]
e os núcleos se movem então no potencial $E_{\mathrm{el}}(\mathbf R) +
V_{\mathrm{nn}}(\mathbf R)$, a energia eletrônica mais a repulsão entre os
núcleos.
\end{definition}

\begin{definition}[Superfície de energia potencial, ponto estacionário, otimização de geometria]\label{def:b3:computational-chemistry:pes}
A \emph{superfície de energia potencial}\index{superficie de energia potencial@superfície de energia potencial} (SEP) de
uma molécula é a função $U(\mathbf R) = E_{\mathrm{el}}(\mathbf R) +
V_{\mathrm{nn}}(\mathbf R)$ de suas coordenadas nucleares. Um \emph{ponto
estacionário}\index{ponto estacionario@ponto estacionário} é um ponto em que todas as derivadas primeiras de $U$
se anulam; um mínimo é uma \emph{geometria de equilíbrio}\index{geometria de equilibrio@geometria de equilíbrio}.
Procurar um mínimo a partir de uma estrutura inicial, seguindo as
forças $-\nabla U$, é uma \emph{otimização de geometria}\index{otimizacao de geometria@otimização de geometria}.
\end{definition}

Para uma molécula diatômica, a SEP é uma curva $U(R)$; para uma triatômica,
uma função de três coordenadas; para $N$ átomos, de $3N - 6$. A massa dos
núcleos não entra em $\hat H_{\mathrm{el}}$: \ce{H2}, \ce{HD} e \ce{D2} compartilham
a mesma SEP, o mesmo comprimento de ligação e a mesma \omterm{def:b3:quantum-model-systems:oscillator}{constante de força}, e
diferem apenas no modo como os núcleos se movem sobre ela (\omterm{def:b3:quantum-model-systems:oscillator}{energia do ponto
zero}, \cref{ch:b3:quantum-model-systems}).

\begin{proposition}[Frequências a partir da hessiana]\label{prop:b3:computational-chemistry:hessian}
Perto de um mínimo, $U \approx U_0 + \frac12\sum_{ij}H_{ij}\,\delta q_i\,\delta q_j$, com
a hessiana $H_{ij} = \partial^2U/\partial q_i\partial q_j$. Em coordenadas
ponderadas pelas massas, $\delta q_i\sqrt{m_i}$, os \omterm{def:b3:quantum-model-systems:operator}{autovalores} $\lambda_k$ da
hessiana ponderada pelas massas dão as frequências angulares harmônicas $\omega_k =
\sqrt{\lambda_k}$ dos modos normais. Em um mínimo, todos os $\lambda_k$ são
positivos (seis, ou cinco para uma molécula linear, são nulos: translações e
rotações); em um ponto de sela de primeira ordem, exatamente um é negativo.
\end{proposition}

\begin{proof}[Demonstração parcial]
Para uma diatômica, $U \approx U_0 + \frac12k(R - R_e)^2$ com $k = U''(R_e)$, e
o oscilador do \cref{ch:b3:quantum-model-systems} dá $\omega =
\sqrt{k/\mu}$. Em geral, as equações clássicas $m_i\ddot q_i =
-\sum_jH_{ij}q_j$ tornam-se, nas coordenadas ponderadas pelas massas $x_i = \sqrt{m_i}q_i$,
$\ddot{\mathbf x} = -\tilde H\mathbf x$ com a matriz simétrica $\tilde H_{ij} =
H_{ij}/\sqrt{m_im_j}$; seus autovetores oscilam independentemente a
$\sqrt{\lambda_k}$. Um $\lambda_k$ negativo dá uma frequência imaginária: nessa
direção a energia diminui, é um máximo. Os modos normais do
\cref{ch:b3:group-theory-applied} são esses autovetores.
\end{proof}

\begin{omfigure}
\begin{tikzpicture}[font=\small, scale=0.95]
  \fill[omDef!8] (0,0) ellipse (1.9 and 1.15);
  \foreach \s in {1.0,0.75,0.5,0.25}{\draw[omDef!70] (0,0) ellipse ({1.9*\s} and {1.15*\s});}
  \fill[omThm!8] (6,0.4) ellipse (1.6 and 1.0);
  \foreach \s in {1.0,0.75,0.5,0.25}{\draw[omThm!70] (6,0.4) ellipse ({1.6*\s} and {1.0*\s});}
  \draw[gray] (2.1,-0.15) .. controls (2.9,0.55) and (3.4,0.55) .. (4.3,0.35);
  \draw[gray] (2.2,0.9) .. controls (3.0,0.35) and (3.5,0.35) .. (4.1,1.2);
  \node[omsaddle] at (3.15,0.42) {};
  \draw[ompath] (0,0) .. controls (1.8,0.6) and (2.4,0.5) .. (3.15,0.42) .. controls (4.0,0.35) and (4.8,0.4) .. (6,0.4);
  \node[anchor=north] at (0,-1.25) {mínimo A (reagentes)};
  \node[anchor=north] at (6,-0.7) {mínimo B (produtos)};
  \node[anchor=south] at (3.15,1.25) {ponto de sela};
  \draw[->, gray] (3.15,1.25) -- (3.15,0.6);
  \fill (0,0) circle (1.2pt) (6,0.4) circle (1.2pt);
  \draw[->] (-2.3,-2.2) -- (8.2,-2.2) node[right] {$q_1$};
  \draw[->] (-2.3,-2.2) -- (-2.3,1.8) node[above] {$q_2$};
\end{tikzpicture}

{\small Uma \omterm{def:b3:computational-chemistry:pes}{superfície de energia potencial} vista de cima, como um mapa de
curvas de nível sobre duas coordenadas nucleares (esquema). Dois mínimos são
ligados, através de um ponto de sela (cruz), pelo caminho de menor energia
(tracejado): o caminho de energia mínima estudado no \cref{ch:b3:rate-theories}.}
\end{omfigure}

\section{O princípio variacional}

A \omterm{thm:b3:quantum-model-systems:schrodinger}{equação de Schrödinger} eletrônica não pode ser resolvida exatamente para
nenhuma molécula com mais de um elétron. As aproximações são julgadas por um
único teorema.

\begin{theorem}[O princípio variacional]\label{thm:b3:computational-chemistry:variational}
Para toda função tentativa normalizável $\phi$ que satisfaz as condições de
contorno do problema,
\[
E[\phi] = \frac{\langle\phi|\hat H|\phi\rangle}{\langle\phi|\phi\rangle} \ge E_0,
\]
onde $E_0$ é o menor \omterm{def:b3:quantum-model-systems:operator}{autovalor} de $\hat H$; a igualdade só vale se $\phi$
for uma \omterm{def:b3:quantum-model-systems:operator}{autofunção} do estado fundamental. É o \emph{princípio
variacional}\index{principio variacional@princípio variacional}.
\end{theorem}

\begin{proof}
Expanda $\phi = \sum_nc_n\psi_n$ nas \omterm{def:b3:quantum-model-systems:operator}{autofunções} ortonormais de $\hat H$
(\cref{thm:b3:quantum-model-systems:hermitian-real}). Então $\langle\phi|\hat
H|\phi\rangle = \sum_n|c_n|^2E_n \ge E_0\sum_n|c_n|^2 = E_0\langle\phi|\phi\rangle$,
pois todo $E_n \ge E_0$; a igualdade exige $c_n = 0$ sempre que $E_n > E_0$.
\end{proof}

Quanto mais baixa a energia de uma função tentativa, melhor ela é; todo
parâmetro ajustável é fixado minimizando $E$.

\begin{proposition}[O hélio com uma carga blindada]\label{prop:b3:computational-chemistry:helium-zeta}
Para o hélio, a função tentativa $\phi = \eu^{-\zeta(r_1 + r_2)}$ (em unidades atômicas:
comprimentos em $a_0$, energias em $E_h$) dá $E(\zeta) = \zeta^2 - \frac{27}{8}\zeta$,
mínima para $\zeta = \frac{27}{16}$, com $E = -(\frac{27}{16})^2 E_h =
\qty{-2.8477}{\hartree}$.
\end{proposition}

\begin{proof}
Para um elétron em $\eu^{-\zeta r}$, $\langle T\rangle = \zeta^2/2$ e
$\langle 1/r\rangle = \zeta$; a repulsão de dois elétrons na mesma função $1s$
vale $\frac58\zeta$ (uma integral clássica, admitida). Logo $E = 2\cdot\frac12
\zeta^2 - 2\cdot2\zeta + \frac58\zeta = \zeta^2 - \frac{27}{8}\zeta$, e $\dd E/\dd
\zeta = 0$ dá $\zeta = 27/16 = 1.6875$.
\end{proof}

\begin{example}[Quão boa ela é?]\label{ex:b3:computational-chemistry:helium}
% ledger: ie:He, ie4:He2, const:Eh, const:eV
A energia total do hélio é o oposto da soma de suas duas energias de ionização,
$-(24.5874 + 54.4178)\,\unit{eV} = \qty{-2.9034}{\hartree}$. A função de
um parâmetro fica \qty{0.056}{\hartree} (\qty{1.5}{eV}) alta demais: cada
elétron vê um núcleo blindado pelo outro até $Z_{\mathrm{eff}} = 1.69$, como
sugerem as regras de Slater, mas os elétrons também se esquivam um do outro
de um modo que nenhum produto de duas funções consegue descrever.
\end{example}

\begin{proposition}[Variação linear]\label{prop:b3:computational-chemistry:linear-variation}
Para uma função tentativa $\phi = \sum_{i=1}^n c_i\chi_i$ construída sobre $n$ funções fixas,
com $H_{ij} = \langle\chi_i|\hat H|\chi_j\rangle$ e $S_{ij} = \langle\chi_i|\chi_j
\rangle$, os valores estacionários de $E$ são as raízes de $\det(H - ES) = 0$, e
a menor raiz é um limite superior para $E_0$.
\end{proposition}

\begin{proof}
$E\sum_{ij}c_ic_jS_{ij} = \sum_{ij}c_ic_jH_{ij}$ (coeficientes reais). Derive
em relação a $c_k$ e imponha $\partial E/\partial c_k = 0$: $\sum_j(H_{kj} -
ES_{kj})c_j = 0$ para todo $k$. Uma solução não nula exige um determinante
nulo. A menor raiz é $E[\phi]$ para seu autovetor, logo $\ge E_0$.
\end{proof}

O método de Hückel do volume do 2.º ano é esta proposição com orbitais $p$
e $H_{ij}$ empíricos. Todos os métodos a seguir são a mesma ideia com
funções melhores e elementos de matriz melhores.

\section{Conjuntos de base}

\begin{definition}[Conjunto de base, base mínima, orbitais do tipo Slater e do tipo gaussiano]\label{def:b3:computational-chemistry:basis}
Um \emph{conjunto de base}\index{conjunto de base} é o conjunto das funções fixas $\chi_i$,
centradas nos átomos, a partir das quais se constroem os orbitais
moleculares. Uma \emph{base mínima}\index{base minima@base mínima} tem uma função por
orbital atômico ocupado (um $1s$ para H, cinco para C). Um \emph{orbital do tipo
Slater}\index{orbital do tipo Slater} (STO) tem a forma radial
$r^{n-1}\eu^{-\zeta r}$; um \emph{orbital do tipo gaussiano}\index{orbital do tipo gaussiano}
(GTO) tem a forma $\eu^{-\alpha r^2}$ multiplicada por um polinômio em $x$, $y$,
$z$. Uma \emph{função gaussiana contraída}\index{funcao gaussiana contraida@função gaussiana contraída}
é uma combinação fixa $\sum_kd_k\,\eu^{-\alpha_kr^2}$ de gaussianas primitivas.
\end{definition}

As funções de Slater têm a forma correta (um bico no núcleo, uma cauda
exponencial), mas suas integrais de dois elétrons sobre vários centros são
muito custosas; as funções gaussianas têm a forma errada, mas o produto de
duas gaussianas em centros diferentes é uma gaussiana centrada em um ponto
entre eles, de modo que toda integral se reduz a uma fórmula fechada. As
contrações tiram o melhor dos dois mundos: na base STO-3G, cada função de
Slater é substituída por uma soma fixa de três gaussianas ajustadas a ela.

\begin{example}[A função STO-3G do hidrogênio]\label{ex:b3:computational-chemistry:sto3g}
% ledger: bse:H.sto3g.a1, bse:H.sto3g.a2, bse:H.sto3g.a3, bse:sto3g.c1, bse:sto3g.c2, bse:sto3g.c3
A função $1s$ do hidrogênio em STO-3G é $0.15433\,g(3.42525) + 0.53533\,
g(0.62391) + 0.44463\,g(0.16886)$, onde $g(\alpha)$ é a gaussiana normalizada
$(2\alpha/\pi)^{3/4}\eu^{-\alpha r^2}$ (expoentes em $a_0^{-2}$). Ela imita uma
função de Slater de expoente $\zeta = 1.24$ --- um átomo de hidrogênio
ligeiramente comprimido, como ele está nas moléculas. Sua sobreposição com a
função $1s$ exata de mesmo expoente é $0.9998$; só o bico no núcleo e a cauda
distante ficam de fora (figura abaixo).
\end{example}

\begin{definition}[Valência dividida, funções de polarização e funções difusas]\label{def:b3:computational-chemistry:split-valence}
Uma \emph{base de valência dividida}\index{base de valencia dividida@base de valência dividida} descreve cada
orbital de valência por duas (ou mais) funções de tamanhos diferentes, de
modo que o orbital possa crescer ou encolher em uma molécula. Uma \emph{função de
polarização}\index{funcao de polarizacao@função de polarização} tem um $l$ maior que o dos orbitais
ocupados ($p$ em H, $d$ em C) e permite que a densidade se desloque para fora do átomo; uma
\emph{função difusa}\index{funcao difusa@função difusa} tem um expoente pequeno e
descreve ânions e interações fracas longe dos núcleos.
\end{definition}

Um nome como 6-31G(d) diz: seis primitivas para cada função de caroço, a
valência dividida em uma contração de três e uma única primitiva, e \omterm{def:b3:computational-chemistry:split-valence}{funções
de polarização} $d$ nos átomos pesados. Os resultados convergem para o
\emph{limite de base completa} à medida que a base cresce; um cálculo só é
tão bom quanto sua base.

\begin{omfigure}
\begin{tikzpicture}
\begin{axis}[width=0.46\linewidth, height=5.4cm, xmin=0, xmax=5, ymin=0,
  ymax=0.6, axis lines=left, xlabel={$r$ / $a_0$}, ylabel={$\phi(r)$},
  legend cell align=left, legend style={font=\small, draw=none},
  tick label style={font=\small}, label style={font=\small},
  title={\small STO-3G comparada com $\eu^{-r}$ ($\zeta = 1$)}]
\addplot[black, thick] table[x=r, y=slater] {figdata/out/bachelor-3-minimal-basis-hf-sto.dat};
\addlegendentry{Slater}
\addplot[omDef, thick, dashed] table[x=r, y=g3] {figdata/out/bachelor-3-minimal-basis-hf-sto.dat};
\addlegendentry{STO-3G}
\addplot[omThm, thick, dotted] table[x=r, y=g1] {figdata/out/bachelor-3-minimal-basis-hf-sto.dat};
\addlegendentry{uma gaussiana}
\end{axis}
\end{tikzpicture}\hfill
\begin{tikzpicture}
\begin{axis}[width=0.46\linewidth, height=5.4cm, xmin=0.5, xmax=8.5,
  ymin=-2.86, ymax=-2.78, axis lines=left, xlabel={ciclo SCF},
  ylabel={$E$ / $E_h$}, ytick={-2.85,-2.83,-2.81,-2.79},
  yticklabel style={/pgf/number format/fixed, /pgf/number format/precision=2},
  tick label style={font=\small}, label style={font=\small},
  title={\small SCF do \ce{HeH+}}]
\addplot[omDef, thick, mark=*] table[x=it, y=E] {figdata/out/bachelor-3-minimal-basis-hf-heh-scf.dat};
\addplot[gray, dashed, domain=0.5:8.5] {-2.84184};
\end{axis}
\end{tikzpicture}

{\small À esquerda: a contração STO-3G (tracejada) reproduz de perto a função
de Slater $1s$, exceto no núcleo; uma única gaussiana (pontilhada) não.
À direita: a energia de \ce{HeH+} a \qty{1.4632}{\bohr} durante as iterações
do \omterm{def:b3:computational-chemistry:hartree-fock}{campo autoconsistente} (STO-3G); ela converge por cima para
\qty{-2.8418}{\hartree} em poucos ciclos, cada iterado obedecendo ao \omterm{thm:b3:computational-chemistry:variational}{princípio
variacional}.}
\end{omfigure}

\section{A teoria de Hartree--Fock}

A \omterm{def:b3:many-electron-atoms:indistinguishable}{função de onda antissimétrica} mais simples de $N$ elétrons é um único
\omterm{def:b3:many-electron-atoms:slater}{determinante de Slater}. A teoria de Hartree--Fock encontra o melhor deles.

\begin{definition}[Método de Hartree--Fock, operador de Fock, campo autoconsistente]\label{def:b3:computational-chemistry:hartree-fock}
O \emph{método de Hartree--Fock}\index{metodo de Hartree--Fock@método de Hartree--Fock} aproxima o
estado fundamental pelo \omterm{def:b3:many-electron-atoms:slater}{determinante de Slater} único de menor energia. Seus
orbitais são \omterm{def:b3:quantum-model-systems:operator}{autofunções} do \emph{operador de Fock}\index{operador de Fock}
\[
\hat f(1) = \hat h(1) + \sum_b\big[2\hat J_b(1) - \hat K_b(1)\big],
\]
no caso de camadas fechadas: $\hat h$ é a energia cinética e a atração pelos
núcleos de um elétron, $\hat J_b$ a repulsão pela nuvem de carga do
orbital $b$, e $\hat K_b$ o \omterm{def:b3:quantum-model-systems:operator}{operador} de troca, cujos \omterm{def:b3:quantum-model-systems:hermitian}{valores esperados} são
as \omterm{def:b3:many-electron-atoms:exchange}{integrais de troca} do \cref{ch:b3:many-electron-atoms}. Como $\hat f$
depende dos orbitais que ele determina, as equações são resolvidas por
iteração até que os orbitais não mudem mais: um \emph{campo
autoconsistente}\index{campo autoconsistente} (SCF).
\end{definition}

\begin{theorem}[Equações de Roothaan--Hall]\label{thm:b3:computational-chemistry:roothaan}
Escrevendo cada orbital em uma base, $\phi_a = \sum_\mu C_{\mu a}\chi_\mu$, as
equações de Hartree--Fock de camadas fechadas tornam-se o problema matricial de \omterm{def:b3:quantum-model-systems:operator}{autovalores}
\[
\mathbf F\mathbf C = \mathbf S\mathbf C\boldsymbol\varepsilon, \qquad F_{\mu\nu} =
h_{\mu\nu} + \sum_{\lambda\sigma}P_{\lambda\sigma}\big[(\mu\nu|\sigma\lambda) -
\tfrac12(\mu\lambda|\sigma\nu)\big],
\]
com a matriz densidade $P_{\lambda\sigma} = 2\sum_a^{\mathrm{occ}}C_{\lambda a}C_{\sigma
a}$ e as integrais de dois elétrons $(\mu\nu|\lambda\sigma) = \iint\chi_\mu(1)\chi_\nu(1)
r_{12}^{-1}\chi_\lambda(2)\chi_\sigma(2)$ (unidades atômicas).
\end{theorem}

\begin{proof}[Demonstração parcial]
A energia do determinante é $E = \sum_{\mu\nu}P_{\mu\nu}h_{\mu\nu} + \frac12
\sum P_{\mu\nu}P_{\lambda\sigma}[(\mu\nu|\sigma\lambda) - \frac12(\mu\lambda|\sigma\nu)]$, uma
função quadrática dos coeficientes. Minimize-a sob as restrições de que os
orbitais permaneçam ortonormais, $\mathbf C^{\mathsf T}\mathbf S\mathbf C = \mathbf 1$,
com um multiplicador de Lagrange para cada restrição: as condições de
estacionariedade são $\mathbf F\mathbf C = \mathbf S\mathbf C\boldsymbol\varepsilon$, e
os multiplicadores formam uma matriz que pode ser diagonalizada por uma rotação
dos orbitais ocupados entre si. A álgebra é admitida.
\end{proof}

\begin{method}[O procedimento SCF]\label{met:b3:computational-chemistry:scf}
\begin{enumerate}
  \item Calcule uma única vez as integrais $S_{\mu\nu}$, $h_{\mu\nu}$ e $(\mu\nu|\lambda\sigma)$.
  \item Faça uma estimativa inicial da matriz densidade (por exemplo, a partir dos
        autovetores de $\mathbf h$ sozinha, a estimativa ``de caroço'').
  \item Construa $\mathbf F$ a partir de $\mathbf P$; resolva $\mathbf F\mathbf C = \mathbf S\mathbf C
        \boldsymbol\varepsilon$ (ortogonalize com $\mathbf S^{-1/2}$, depois diagonalize).
  \item Preencha os orbitais mais baixos, com dois elétrons cada; forme a nova $\mathbf P$
        e a energia.
  \item Repita a partir da etapa 3 até que a energia e $\mathbf P$ variem menos
        que um limiar.
\end{enumerate}
\end{method}

\begin{omfigure}
\begin{tikzpicture}[font=\small, node distance=4mm,
  box/.style={draw=black!60, rounded corners=2pt, fill=omDef!6, align=center,
  text width=4.4cm, inner sep=4pt}]
  \node[box] (a) {integrais $S$, $h$, $(\mu\nu|\lambda\sigma)$};
  \node[box, below=of a] (b) {estimativa de $\mathbf P$};
  \node[box, below=of b] (c) {construir $\mathbf F(\mathbf P)$};
  \node[box, below=of c] (d) {resolver $\mathbf F\mathbf C = \mathbf S\mathbf C\boldsymbol\varepsilon$};
  \node[box, below=of d] (e) {nova $\mathbf P$, energia $E$};
  \node[diamond, aspect=2.4, draw=black!60, fill=omThm!6, below=of e, inner sep=1pt] (f) {convergiu?};
  \node[box, below=of f, fill=omProp!8] (g) {orbitais, $\varepsilon$, $E$};
  \draw[->] (a) -- (b); \draw[->] (b) -- (c); \draw[->] (c) -- (d);
  \draw[->] (d) -- (e); \draw[->] (e) -- (f);
  \draw[->] (f) -- node[left] {sim} (g);
  \draw[->] (f.east) -- ++(1.6,0) |- node[pos=0.25, right] {não} (c.east);
\end{tikzpicture}

{\small O ciclo do \omterm{def:b3:computational-chemistry:hartree-fock}{campo autoconsistente} de um cálculo Hartree--Fock.}
\end{omfigure}

\begin{example}[\ce{H2} na base mínima]\label{ex:b3:computational-chemistry:h2}
% ledger: re:H2, janaf:H.dfH0, diat:H2.we, diat:H2.wexe
Com uma função STO-3G por átomo, o orbital ligante é fixado pela simetria,
$\sigma_g \propto \chi_A + \chi_B$, e o SCF precisa de uma única etapa. Em $R =
\qty{1.4}{\bohr}$, a sobreposição é $S_{AB} = 0.6593$ e a energia,
\qty{-1.1167}{\hartree}; o mínimo fica em \qty{1.346}{\bohr} (\qty{71.2}{pm},
contra os \qty{74.1}{pm} medidos) e \qty{-1.1175}{\hartree}. A figura abaixo
compara a curva com a verdadeira, construída a partir de dados espectroscópicos
medidos e da energia de dissociação (uma curva de Morse, \cref{ch:b3:rovibrational-spectroscopy}).
\end{example}

\begin{omfigure}
\begin{tikzpicture}
\begin{axis}[width=0.7\linewidth, height=6cm, xmin=0.8, xmax=10, ymin=-1.2,
  ymax=-0.45, axis lines=left, xlabel={$R$ / $a_0$}, ylabel={$E$ / $E_h$},
  legend cell align=left, legend style={font=\small, at={(0.04,0.98)}, anchor=north west, draw=none},
  tick label style={font=\small}, label style={font=\small}]
\addplot[omDef, thick] table[x=R, y=rhf] {figdata/out/bachelor-3-minimal-basis-hf-h2.dat};
\addlegendentry{RHF/STO-3G}
\addplot[black, thick, dashed] table[x=R, y=morse] {figdata/out/bachelor-3-minimal-basis-hf-h2.dat};
\addlegendentry{medida (forma de Morse)}
\addplot[gray, dotted, domain=0.8:10] {-1.0};
\addplot[omThm, dotted, domain=0.8:10] {-0.93316};
\node[font=\scriptsize, anchor=north] at (axis cs:8,-1.005) {dois átomos de H exatos};
\node[font=\scriptsize, anchor=south, omThm] at (axis cs:8,-0.930) {dois átomos de H STO-3G};
\end{axis}
\end{tikzpicture}

{\small A energia de \ce{H2} em função da distância internuclear: Hartree--Fock
restrito na base STO-3G (traço contínuo) e a curva verdadeira (tracejada,
forma de Morse construída a partir de $D_0$, $\omega_e$, $\omega_ex_e$, $r_e$
medidos). Perto do mínimo o RHF é bom; para $R$ grande, ele sobe muito acima
de dois átomos de hidrogênio.}
\end{omfigure}

\begin{proposition}[Por que o Hartree--Fock restrito falha na dissociação]\label{prop:b3:computational-chemistry:rhf-dissociation}
Na \omterm{def:b3:computational-chemistry:basis}{base mínima}, a função de onda RHF de \ce{H2} é, a menos da
normalização, $\sigma_g(1)\sigma_g(2) \propto \chi_A(1)\chi_B(2) + \chi_B(1)\chi_A(2)
+ \chi_A(1)\chi_A(2) + \chi_B(1)\chi_B(2)$: a qualquer distância, metade covalente e
metade iônica (\ce{H+ H-}). Para $R$ grande, sua energia tende à média da de dois
átomos neutros e da de um par de íons, e não à de dois átomos.
\end{proposition}

\begin{proof}
Desenvolva $(\chi_A + \chi_B)(1)(\chi_A + \chi_B)(2)$: os dois termos cruzados põem
um elétron em cada átomo, os dois termos quadrados põem ambos no mesmo átomo,
com pesos iguais qualquer que seja $R$. O determinante não pode mudar esses pesos.
\end{proof}

\begin{definition}[Correlação eletrônica, energia de correlação]\label{def:b3:computational-chemistry:correlation}
A \emph{correlação eletrônica}\index{correlacao eletronica@correlação eletrônica} é a parte do
afastamento mútuo dos elétrons que um único determinante não descreve (nele,
os elétrons de spins opostos se movem independentemente, cada um no campo
médio do outro). A \emph{energia de correlação}\index{energia de correlacao@energia de correlação} é a energia
exata menos a energia de Hartree--Fock em uma base completa; ela é sempre
negativa.
\end{definition}

A \omterm{def:b3:computational-chemistry:correlation}{energia de correlação} é de cerca de \qty{-0.04}{\hartree} por par de
elétrons, um por cento de uma energia total, mas comparável à energia de uma
reação: os métodos que a incluem (interação de configurações, teoria de
perturbação, \textit{coupled cluster}) custam muito mais que o Hartree--Fock, cujo
custo já cresce aproximadamente com a quarta potência do número de funções de base.

\begin{theorem}[Teorema de Koopmans]\label{thm:b3:computational-chemistry:koopmans}
Na teoria de Hartree--Fock, a energia necessária para retirar um elétron do
orbital ocupado $a$, mantendo todos os outros orbitais congelados, é $-\varepsilon_a$. As
energias de ionização medidas são, portanto, aproximadamente os opostos das
energias orbitais: é o \emph{teorema de Koopmans}\index{teorema de Koopmans}.
\end{theorem}

\begin{proof}
A energia de um determinante de camadas fechadas é $\sum_a2h_{aa} + \sum_{ab}(2J_{ab} -
K_{ab})$, e $\varepsilon_a = h_{aa} + \sum_b(2J_{ab} - K_{ab})$. Retire um elétron
de $a$ sem mudar os orbitais: os termos perdidos são $h_{aa}$, as
interações desse elétron com todos os outros, $\sum_b(2J_{ab} - K_{ab}) -
J_{aa}$, mais $J_{aa}$ com seu antigo parceiro --- no total, exatamente
$\varepsilon_a$. Logo $E^+ - E = -\varepsilon_a$.
\end{proof}

\begin{example}[Ionizando \ce{H2}]\label{ex:b3:computational-chemistry:koopmans}
% ledger: ie4:H2
A \qty{1.4}{\bohr}, a energia orbital STO-3G de $\sigma_g$ é
\qty{-0.5782}{\hartree}: Koopmans prevê \qty{15.73}{eV}, contra os
\qty{15.43}{eV} medidos. Os orbitais congelados superestimam a energia do íon
(ele relaxaria), e a correlação ausente age no sentido oposto; para as
ionizações de valência os dois erros se compensam em grande parte, e é por
isso que os espectros de fotoelétrons do volume do 2.º ano podem ser lidos com
diagramas de orbitais.
\end{example}

\section{A teoria do funcional da densidade, e o que um cálculo pode dizer}

A função de onda de $N$ elétrons depende de $3N$ coordenadas; a \omterm{def:b3:computational-chemistry:dft}{densidade
eletrônica} depende de três.

\begin{definition}[Densidade eletrônica, teoria do funcional da densidade]\label{def:b3:computational-chemistry:dft}
A \emph{densidade eletrônica}\index{densidade eletronica@densidade eletrônica} $\rho(\mathbf r)$ é o
número de elétrons por unidade de volume em $\mathbf r$, somado sobre todos os elétrons:
$\int\rho\,\dd\tau = N$. A \emph{teoria do funcional da densidade}\index{teoria do funcional da densidade}
(DFT) calcula a energia do estado fundamental como um funcional de $\rho$. Na
prática, $\rho = \sum_a|\phi_a|^2$ é construída a partir dos \emph{orbitais de
Kohn--Sham}\index{orbitais de Kohn--Sham}, os orbitais de elétrons independentes
fictícios que têm a mesma densidade que os reais, e todos os efeitos de troca
e de correlação são reunidos em um \emph{funcional de
troca--correlação}\index{funcional de troca--correlacao@funcional de troca--correlação} $E_{\mathrm{xc}}[\rho]$.
\end{definition}

\begin{theorem}[Hohenberg--Kohn]\label{thm:b3:computational-chemistry:hohenberg-kohn}
A densidade do estado fundamental de um sistema de elétrons determina o
potencial externo (os núcleos) a menos de uma constante, logo o hamiltoniano e
todas as propriedades do estado fundamental; e o funcional de energia da
densidade é mínimo na densidade verdadeira.
\end{theorem}
\admitted

A demonstração é tratada em cursos mais avançados. Ela garante que existe um
funcional exato, mas não diz qual é: todo funcional prático é uma
aproximação, escolhida em uma hierarquia --- densidade local, correções de
gradiente, híbridos que incorporam uma parte de troca Hartree--Fock --- e
calibrada sobre dados de referência. As equações de Kohn--Sham são resolvidas
pelo mesmo ciclo SCF que o Hartree--Fock, a um custo semelhante, mas incluem a
correlação; é por isso que a DFT se tornou o cavalo de batalha da química.

\begin{method}[Ler um cálculo]\label{met:b3:computational-chemistry:read-calculation}
\begin{enumerate}
  \item Verifique que o SCF convergiu e que a \omterm{def:b3:computational-chemistry:pes}{otimização de geometria} terminou
        com forças abaixo do limiar.
  \item Faça um cálculo de frequências no mesmo nível: todas as frequências
        reais indicam um mínimo; uma frequência imaginária, uma estrutura de
        transição, cujo modo normal é o movimento que atravessa a barreira.
  \item Compare o que é comparável: diferenças de energia entre estruturas
        calculadas em um mesmo nível, nunca energias absolutas de níveis diferentes.
  \item Compare com o experimento conhecendo os erros típicos do nível:
        as frequências harmônicas costumam ser multiplicadas por um fator de 0.9 a 0.97;
        os comprimentos de ligação são bons a cerca de 1 pm com um funcional híbrido e
        uma base polarizada; as energias de reação, a alguns \unit{kJ/mol}, no melhor caso.
\end{enumerate}
\end{method}

\begin{method}[Escolher um nível de teoria]\label{met:b3:computational-chemistry:choose-level}
Geometrias e frequências de moléculas orgânicas: um funcional híbrido com uma
\omterm{def:b3:computational-chemistry:split-valence}{base de valência dividida} polarizada. Interações fracas: acrescente uma
correção de dispersão e \omterm{def:b3:computational-chemistry:split-valence}{funções difusas}. Barreiras de reação e energias
precisas: métodos correlacionados de função de onda sobre geometrias DFT,
até onde o orçamento permitir. Sistemas grandes (proteínas, superfícies): uma
pequena região tratada pela mecânica quântica dentro de um campo de forças
clássico. Teste sempre o nível escolhido em uma molécula aparentada cuja
resposta é conhecida.
\end{method}

\begin{inthelab}[Um experimento computacional]
Um cálculo é conduzido como um experimento e registrado no caderno: o
programa e sua versão, o método, a base, os limiares de convergência, a
estrutura inicial, e os arquivos de saída guardados como dados brutos. Uma
sequência típica é: construir a molécula, otimizar sua geometria em um nível
modesto, calcular as frequências para confirmar um mínimo e, por fim,
calcular a energia em um nível mais alto sobre essa geometria.
\end{inthelab}

\begin{history}[Uma molécula conhecida primeiro pelo cálculo]
% ledger: hist4:HeHp.2019
\ce{HeH+} foi produzido em laboratório em 1925. Suas propriedades foram
calculadas em detalhe a partir da década de 1930, e previu-se que seria
abundante em certos gases astrofísicos; mas sua primeira raia rotacional fica
no infravermelho distante, absorvida pela atmosfera. Em 2019, um
observatório aerotransportado o detectou na nebulosa planetária NGC 7027, na
frequência que os cálculos e os espectros de laboratório haviam previsto.
\end{history}

\section{Exercícios}

\begin{exercise}[$\star$]\label{exo:b3:computational-chemistry:1}
Para o átomo de hidrogênio, tome a função tentativa $\eu^{-\alpha r^2}$. Em
unidades atômicas, $E(\alpha) = \frac32\alpha - 2\sqrt{2\alpha/\pi}$. Encontre o melhor
$\alpha$ e a energia correspondente; compare com $-\frac12E_h$.
\end{exercise}

\begin{exercise}[$\star$]\label{exo:b3:computational-chemistry:2}
Quantas funções de base usa um cálculo sobre o benzeno na base STO-3G?
E em 6-31G(d), com seis funções $d$ cartesianas por carbono e duas funções $s$
por hidrogênio?
\end{exercise}

\begin{exercise}[$\star$]\label{exo:b3:computational-chemistry:3}
Explique por que \ce{H2} e \ce{D2} têm o mesmo comprimento de ligação e a
mesma \omterm{def:b3:quantum-model-systems:oscillator}{constante de força}, mas números de onda vibracionais diferentes. Que
razão você espera?
\end{exercise}

\begin{exercise}[$\star$]\label{exo:b3:computational-chemistry:4}
Um cálculo de frequências sobre uma estrutura de fórmula \ce{C2H5F} dá 17
frequências reais e uma frequência imaginária, $487\iu$~\unit{cm^{-1}}. Que tipo
de \omterm{def:b3:computational-chemistry:pes}{ponto estacionário} é esse? Quantas frequências você esperava no total?
\end{exercise}

\begin{exercise}[$\star\star$]\label{exo:b3:computational-chemistry:5}
% ledger: ie:He, ie4:He2
Usando $E(\zeta) = \zeta^2 - \frac{27}{8}\zeta$, calcule a melhor energia
variacional do hélio em \unit{\hartree} e em eV, e seu erro em relação ao valor exato
$\qty{-2.9034}{\hartree}$. Que fração da energia total representa o erro?
\end{exercise}

\begin{exercise}[$\star\star$]\label{exo:b3:computational-chemistry:6}
Duas funções de base dão $H_{11} = \qty{-1.0}{\hartree}$, $H_{22} =
\qty{-0.5}{\hartree}$, $H_{12} = \qty{-0.2}{\hartree}$, $S_{12} = 0.3$. Resolva a
equação secular. A menor raiz fica abaixo de $H_{11}$?
\end{exercise}

\begin{exercise}[$\star\star$]\label{exo:b3:computational-chemistry:7}
Um cálculo Hartree--Fock leva 10 minutos para o benzeno em STO-3G. Supondo
um custo proporcional à quarta potência do número de funções de base,
estime o tempo em 6-31G(d).
\end{exercise}

\begin{exercise}[$\star\star$]\label{exo:b3:computational-chemistry:8}
% ledger: ie4:H2
A energia orbital STO-3G de \ce{H2} a \qty{1.4}{\bohr} é
\qty{-0.5782}{\hartree}. Dê a estimativa de Koopmans da energia de ionização em eV
e compare com \qty{15.43}{eV}. Cite duas razões para a diferença e o sentido
de cada uma.
\end{exercise}

\begin{exercise}[$\star\star$]\label{exo:b3:computational-chemistry:9}
Em STO-3G, a energia de um átomo de hidrogênio é \qty{-0.4666}{\hartree} e a energia
RHF de \ce{H2} a \qty{10}{\bohr} é \qty{-0.5960}{\hartree}. Quão acima de dois
átomos fica a energia RHF a essa distância, em eV? Explique com os termos
iônicos.
\end{exercise}

\begin{exercise}[$\star\star\star$]\label{exo:b3:computational-chemistry:10}
Para $\phi = c_1\chi_1 + c_2\chi_2$ com funções reais, escreva $E(c_1,c_2)$ e
deduza as duas equações seculares impondo $\partial E/\partial c_1 =
\partial E/\partial c_2 = 0$.
\end{exercise}

\begin{exercise}[$\star\star\star$]\label{exo:b3:computational-chemistry:11}
% ledger: aw:H, diat:H2.we
A energia RHF/STO-3G de \ce{H2} é $-1.116871$, $-1.117501$ e
$\qty{-1.116714}{\hartree}$ em $R = 1.30$, 1.35 e \qty{1.40}{\bohr}. Estime
a \omterm{def:b3:quantum-model-systems:oscillator}{constante de força} por uma diferença finita, converta-a em \unit{N/m}
($\qty{1}{\hartree\per\bohr\squared} = \qty{1556.9}{N/m}$) e calcule o
número de onda harmônico. Compare com $\tilde\omega_e = \qty{4401}{cm^{-1}}$.
\end{exercise}

\begin{exercise}[$\star\star\star$]\label{exo:b3:computational-chemistry:12}
Você precisa calcular a diferença de energia entre dois confôrmeros de um
açúcar, que se sabe diferirem de alguns \unit{kJ/mol}, com uma ligação de
hidrogênio intramolecular em um deles. Escolha e justifique um método, um
\omterm{def:b3:computational-chemistry:basis}{conjunto de base} e as verificações que você faria.
\end{exercise}

\section{Problema: HeH\texorpdfstring{$^+$}{+}, a primeira molécula}

\begin{problem}[{Problema de fim de semana --- um cálculo Hartree--Fock em base
mínima do íon hidreto de hélio: a base, a estimativa de caroço, o orbital
convergido e a ligação de um próton ao hélio}]
\label{pb:b3:computational-chemistry:1}
% ledger: bse:He.sto3g.a1, bse:He.sto3g.a2, bse:He.sto3g.a3, bse:H.sto3g.a1, pa:He, ie:He, ie4:He2
\ce{HeH+} é calculado em $R = \qty{1.4632}{\bohr}$ na base STO-3G (expoentes
do He 6.3624, 1.1589, 0.31365; expoentes do H 3.4253, 0.62391, 0.16886; os
mesmos três coeficientes). A função 1 está no He, a função 2 no H. O programa
fornece (unidades atômicas)
\[
\mathbf S = \begin{pmatrix}1 & 0.5368\\ 0.5368 & 1\end{pmatrix},\quad
\mathbf h = \begin{pmatrix}-2.5983 & -1.4318\\ -1.4318 & -1.7318\end{pmatrix},\quad
\mathbf F_{\mathrm{final}} = \begin{pmatrix}-1.5902 & -1.0610\\ -1.0610 & -0.8340
\end{pmatrix},
\]
o orbital ocupado $\phi = 0.8766\chi_1 + 0.2025\chi_2$, as energias orbitais
$-1.6328$ e \qty{-0.1725}{\hartree}, e as energias SCF $-2.7978$, $-2.8404$,
$-2.8418$, $-2.8418$. Um átomo de hélio na mesma base tem
\qty{-2.8078}{\hartree}. $\qty{1}{\hartree} = \qty{27.211}{eV} =
\qty{2625.5}{kJ/mol}$.

\medskip\noindent\textbf{Parte I --- A base.}
\begin{enumerate}
  \item O que significa ``STO-3G''?
  \item Por que os expoentes do hélio são maiores que os do hidrogênio? Verifique
        que sua razão é $(1.69/1.24)^2$.
  \item Quantas funções de base, elétrons e orbitais ocupados há?
  \item Calcule a energia de repulsão nuclear.
  \item Interprete $S_{12} = 0.5368$.
  \item Por que $h_{11}$ é mais baixo que $h_{22}$?
\end{enumerate}

\medskip\noindent\textbf{Parte II --- A estimativa de caroço.}
\begin{enumerate}[resume]
  \item Escreva a equação secular $\det(\mathbf h - \varepsilon\mathbf S) = 0$ como
        uma equação do segundo grau em $\varepsilon$.
  \item Resolva-a.
  \item Por que isso é apenas um ponto de partida?
  \item Dê a fórmula da matriz de Fock em função de $\mathbf h$, da
        matriz densidade e das integrais de dois elétrons.
  \item Que interações entre elétrons $\mathbf F$ contém e que faltam em $\mathbf h$?
  \item Explique por que o procedimento precisa ser iterado.
\end{enumerate}

\medskip\noindent\textbf{Parte III --- O orbital convergido.}
\begin{enumerate}[resume]
  \item Verifique que $\phi$ está normalizado.
  \item As populações são os elementos diagonais de $P\mathbf S$, $2c_1^2 + 2c_1c_2S$
        no He e $2c_2^2 + 2c_1c_2S$ no H. Calcule-as.
  \item Onde fica a carga positiva? A molécula é mais bem descrita
        como \ce{He + H+} ou como \ce{He+ + H}?
  \item Use o \omterm{thm:b3:computational-chemistry:koopmans}{teorema de Koopmans} para estimar a energia necessária para retirar
        um elétron de \ce{HeH+}.
  \item Calcule a energia eletrônica $E_{\mathrm{el}} = E - V_{\mathrm{nn}}$.
  \item Quantos ciclos foram necessários para convergir a \qty{e-4}{\hartree}? Por que
        a energia diminuiu a cada ciclo?
  \item Com os expoentes do hélio reescalados para $\zeta = 2.0925$ em vez de 1.69, o
        mesmo programa dá \qty{-2.8607}{\hartree}. Qual base é melhor, e
        por que se pode afirmá-lo?
  \item O que representa a energia do orbital virtual, \qty{-0.1725}{\hartree}?
\end{enumerate}

\medskip\noindent\textbf{Parte IV --- A ligação de um próton.}
\begin{enumerate}[resume]
  \item Qual é a energia de um próton isolado? E a de \ce{He + H+} muito
        afastados, nesta base?
  \item Calcule a energia liberada por \ce{He + H+ -> HeH+} nesse nível, em
        \unit{kJ/mol}.
  \item A afinidade protônica medida do hélio é \qty{177.8}{kJ/mol}.
        Comente.
  \item A energia exata do hélio é o oposto da soma de suas energias de
        ionização, 24.5874 e \qty{54.4178}{eV}. Calcule-a em \unit{\hartree} e
        o erro do átomo de hélio STO-3G.
  \item Por que um erro desse tamanho ainda pode dar geometrias moleculares
        úteis?
  \item Enuncie o resultado: a energia total RHF/STO-3G de \ce{HeH+} a
        \qty{1.4632}{\bohr}.
\end{enumerate}
\end{problem}
